% Acceptance test for the institutional-research theme + equity-research
% publication type (Steps 2-3 of
% docs/superpowers/specs/2026-09-09-reportkit-institutional-theme-and-equity-profile-spec.md).
% Deliberately compact -- a fast lualatex smoke test that exercises every
% primitive at least once, distinct from
% latex_templates/examples/equity-research/ (the full four-page fictional
% publication, which is the visual-regression fixture, not the acceptance
% smoke test). Requires lualatex; scripts/acceptance_check.sh's
% pdflatex-only TEST_TEXES list must never include this file -- the
% institutional-research theme's own \ifPDFTeX guard would fail it
% immediately under pdflatex, which is the correct behavior (spec open
% question 1), not a bug in this fixture.
\documentclass[theme=institutional-research,publication-type=equity-research]{reportkit}

% Keep the checked-in smoke fixture self-contained when the repository's
% licensed Google Sans fixtures are staged beside it. Consumer projects can
% omit this directory and exercise the theme's configured fallback policy.
\IfFileExists{font_data/GoogleSans-Regular.ttf}{%
  \setreportkitfontpath{font_data/}%
}{}

\setreportkitleftheader{REPORTKIT / INSTITUTIONAL-EQUITY ACCEPTANCE TEST}
\setreportkitfooter{Institutional-research theme + equity-research acceptance test}
\setreportkitversion{v1.8.0}

\title{Institutional-research theme acceptance test}
\author{Test build}
\date{9 September 2026}

\begin{document}

\begin{researchfrontpage}
\researchkicker{Acceptance Testing / Fixture}
\researchheadline{Every equity-research primitive renders once, in one compact document}
\researchdeck{This document is not meant to be read as research -- it exists to give
lualatex one document that touches every primitive Steps 2 and 3 shipped, so a
compile failure here points straight at the theme or publication-type file, not
at a four-page fixture.}

\begin{ratingstrip}
\ratingitem{Stock Rating}{\color{Accent}Overweight}
\ratingitem{Industry View}{Attractive}
\ratingitem{Price Target}{\$100}
\end{ratingstrip}

\begin{whatschanged}
\change{Test Metric}{Old}{New}
\end{whatschanged}

\begin{researchmain}
Main-column body text uses the institutional theme's 10.7pt/14.6pt normal size,
never the sidebar's smaller \texttt{\textbackslash rksidebarsize}.

\begin{principle}{Quieter callouts}
Semantic callouts (spec section 21) restyle quietly under this theme: a thin
left rule and a small uppercase label, not a colored box.
\end{principle}
\end{researchmain}\begin{researchsidebar}
\begin{analystblock}
\sidebarvalue{Test Analyst}\\
\textcolor{Muted}{Fixture Author}
\end{analystblock}

\begin{marketdatablock}
\sidebarrow{Share Price}{\$74.00}
\sidebarrow{Market Cap (\$M)}{1,000}
\end{marketdatablock}

\begin{estimatesblock}
\begin{tabularx}{\linewidth}{@{}Yr@{}}
& \textbf{FY1E}\\
\midrule
EPS (\$) & 1.00\\
\end{tabularx}
\end{estimatesblock}
\end{researchsidebar}
\end{researchfrontpage}

\newpage
\section{Exhibit compositions}

\begin{fullwidthexhibit}
\begin{exhibit}[title={A full-width exhibit contains a financial table}, source={Fixture data.}, label={tab:full}]
\begin{financialtable}
\begin{tabularx}{\linewidth}{@{}Yrr@{}}
\toprule
\textbf{Metric} & \textbf{FY1E} & \textbf{FY2E}\\
\midrule
Revenue (\$M) & 100 & 120\\
\bottomrule
\end{tabularx}
\end{financialtable}
\end{exhibit}
\end{fullwidthexhibit}

\begin{exhibitpair}
\exhibitpane{
\begin{exhibit}[title={Left pane of an exhibit pair}, source={Fixture data.}]
Left pane content.
\end{exhibit}
}
\exhibitpane{
\begin{exhibit}[title={Right pane of an exhibit pair}, source={Fixture data.}]
Right pane content.
\end{exhibit}
}
\end{exhibitpair}

\begin{exhibitgrid}[columns=3]
\exhibitpane{\begin{exhibit}[title={Grid pane one}]Pane one.\end{exhibit}}
\exhibitpane{\begin{exhibit}[title={Grid pane two}]Pane two.\end{exhibit}}
\exhibitpane{\begin{exhibit}[title={Grid pane three}]Pane three.\end{exhibit}}
\end{exhibitgrid}

\newpage
\section{Financial model and risk/reward}

\begin{exhibit}[title={Dense model-page type}, source={Fixture data.}]
\begin{financialmodelpage}
\begin{tabular}{@{}lrr@{}}
\toprule
\textbf{\$M} & \textbf{1Q} & \textbf{2Q}\\
\midrule
Revenue & 25 & 27\\
\bottomrule
\end{tabular}
\end{financialmodelpage}
\end{exhibit}

\rksubheading{Scenarios}
\bullcase{\$130}{Bull case body text.}
\basecase{\$100}{Base case body text.}
\bearcase{\$70}{Bear case body text.}

\end{document}
